%% kapa: exam=1 round=25 year=2014 period=2 form=A subject=appraisal_law subject_ko=부동산관계법규 seq=22 number=22
%% source: Q-Net 감정평가사 기출문제 (cst00309) artlSeq=5101207 file=extracted/1008/5101207/2014년 제25회 감정평가사 1차시험 기출문제/2교시(A형).hwp
%% conversion: kordoc 4.24.51 -> tools/gen1.py (auto); review=proofread-sonnet; flags=none
\begin{question}{22}
부동산 가격공시 및 감정평가에 관한 법령상 개별공시지가의 결정ㆍ공시 등에 관한 설명으로 옳지 않은 것은?
\begin{choices}
  \item 시장ㆍ군수 또는 구청장은 개별공시지가의 조사ㆍ산정지침을 정하여 감정평가업자에게 통보하여야 한다.
  \item 2013년의 공시기준일이 1월 1일인 경우 2013년 5월 15일 토지의 용도변경으로 지목변경이 된 토지에 대한 개별공시지가는 2013년 7월 1일을 기준일로 하여 2013년 10월 31일까지 결정ㆍ공시하여야 한다.
  \item 최근 2년간 업무정지처분을 3회 받은 감정평가업자는 개별토지가격 산정의 타당성 검증을 할 수 없다.
  \item 개별공시지가를 공시하는 시장ㆍ군수 또는 구청장은 필요하다고 인정하는 때에는 개별공시지가의 결정 및 이의신청에 관한 사항을 토지소유자등에게 개별통지할 수 있다.
  \item 감정평가업자는 개별공시지가의 결정을 위한 토지가격의 산정을 위하여 필요한 때에는 타인의 토지에 출입할 수 있다.
\end{choices}
\end{question}
\answer{1}
